% -*- Mode:TeX -*-

%% IMPORTANT: The official thesis specifications are available at:
%%            http://libraries.mit.edu/archives/thesis-specs/
%%
%%            Please verify your thesis' formatting and copyright
%%            assignment before submission.  If you notice any
%%            discrepancies between these templates and the 
%%            MIT Libraries' specs, please let us know
%%            by e-mailing thesis@mit.edu

%% The documentclass options along with the pagestyle can be used to generate
%% a technical report, a draft copy, or a regular thesis.  You may need to
%% re-specify the pagestyle after you \include  cover.tex.  For more
%% information, see the first few lines of mitthesis.cls. 

%\documentclass[12pt,vi,twoside]{mitthesis}
%%
%%  If you want your thesis copyright to you instead of MIT, use the
%%  ``vi'' option, as above.
%%
%\documentclass[12pt,twoside,leftblank]{mitthesis}
%%
%% If you want blank pages before new chapters to be labelled ``This
%% Page Intentionally Left Blank'', use the ``leftblank'' option, as
%% above. 

\documentclass[12pt,twoside]{mitthesis}
\usepackage{lgrind}
%% These have been added at the request of the MIT Libraries, because
%% some PDF conversions mess up the ligatures.  -LB, 1/22/2014
\usepackage{cmap}
\usepackage[T1]{fontenc}

%% Personal packages
\PassOptionsToPackage{american,italian}{babel}
\usepackage{babel}
\selectlanguage{italian}

\usepackage{csquotes}
\PassOptionsToPackage{%
  backend=biber,bibencoding=utf8, %instead of bibtex
  language=auto,%
  style=numeric-comp,%
  %style=authoryear-comp, % Author 1999, 2010
  %bibstyle=authoryear,dashed=false, % dashed: substitute rep. author with ---
  sorting=nyt, % name, year, title
  maxbibnames=10, % default: 3, et al.
  %backref=true,%
  natbib=true % natbib compatibility mode (\citep and \citet still work)
}{biblatex}
\usepackage{biblatex}

\usepackage{amsmath}
\usepackage{amssymb}
\usepackage{ntheorem}
\usepackage{float}
\usepackage{import}
\usepackage{xifthen}
\usepackage{pdfpages}
\usepackage{transparent}

\newcommand{\incfig}[1]{%
    \def\svgwidth{\columnwidth}
    \import{./figures/}{#1.pdf_tex}
}


\let\numberset\mathbb
\newcommand{\N}{\numberset{N}}
\newcommand{\Z}{\numberset{Z}}
\newcommand{\Q}{\numberset{Q}}
\newcommand{\R}{\numberset{R}}
\newcommand{\C}{\numberset{C}}
\newcommand{\U}{\numberset{U}}
\newcommand{\I}{\numberset{I}}
\newcommand{\PP}{\numberset{P}}
\newcommand{\CC}{\complement}
\newcommand{\dom}{\text{dom }}
\newcommand{\im}{\text{Im }}
\newcommand{\foi}{\text{f.i.}}
\newcommand{\e}{\mathrm{e}}
\newcommand{\cc}{\mathcal{C}}
\newcommand{\dd}{\mathrm{d}}

\usepackage{listings}
%\lstset{emph={trueIndex,root},emphstyle=\color{BlueViolet}}%\underbar} % for special keywords
\lstset{language=C++,
  morekeywords={PassOptionsToPackage,selectlanguage},
  basicstyle=\small\ttfamily,
  stringstyle=\rmfamily,
  numbers=none,%left,%
  numberstyle=\scriptsize,%\tiny
  stepnumber=5,
  numbersep=8pt,
  showstringspaces=false,
  breaklines=true,
  %frameround=ftff,
  %frame=single,
  belowcaptionskip=.75\baselineskip
  %frame=L
}
\pagestyle{plain}

\begin{document}
\newtheorem{alg}{Algoritmo}
\newtheorem{theo}{Teorema}
\newtheorem{defn}{Definizione}
\newtheorem{dimo}{Dimostrazione}
\newtheorem{coro}{Corollario}
\newtheorem{prop}{Proposizione}

\title{Algebra lineare e analisi numerica}

\author{Riccardo Cereghino}
\pagestyle{plain}
  % -*- Mode:TeX -*-
%% This file simply contains the commands that actually generate the table of
%% contents and lists of figures and tables.  You can omit any or all of
%% these files by simply taking out the appropriate command.  For more
%% information on these files, see appendix C.3.3 of the LaTeX manual. 
\tableofcontents
%\newpage
%\listoffigures
%\newpage
%\listoftables


\chapter{Accenni al funzionamento di un calcolatore}
Un calcolatore è generalmente composto da:
\begin{itemize}
  \item la $CPU$, esegue le istruzioni, tiene in memoria ($Instruction Register$)
    l'istruzione corrente e la posizione dell'istruzione nella RAM ($P.C.$);
  \item la $RAM$, tiene in memoria le istruzioni da eseguire, certe sezioni
    sono riservate per specifiche funzioni ($stack$, $SP$, $BP$);
  \item dei moduli, come un disco fisso o altro;
  \item la $NIC$, che permette di inviare e ricevere pacchetti in rete.
\end{itemize}

\section{Networking}
Vengono effettuate delle syscall, per eseguire operazioni necessario all'invio
di dati.
Per esempio la syscall send, copia dalla ram una certa sezione di memoria nel
NIC, che verrà inviata in rete.

\subsection{TCP/UDP}
\begin{figure}[t]
\includegraphics[width=.8\textwidth]{gfx/chadudp.jpg}
\centering
\end{figure}

Il protocollo TCP garantisce l'invio del messaggio, dato che attende una
risposta da parte del ricevente.
Il protocollo UDP invia datagrammi, ma non si è certi se la ricezione è
avvenuta.

\subsection{Device driver}
E' un modulo che gestisce le risore fisiche dei vari moduli, tra cui NIC,
gestiscono le interruzioni del dispositivo.

\section{Trasmissione di dati}
\subsection{Invio di dati}
Il DMA presente nel NIC è il componente che accede alla memoria.

La maggior parte dei calcolatori include l'unità di gestione della memoria MMU,
che permette al processore di gestire indirizzi virtuali per ottenere una
maggiore efficienza, ma questo dispositivo è esclusivo del processore, motivo
per cui il DMA del NIC utilizzerà indirizzi fisici.

Delle porzioni di memoria della RAM saranno destinate ad essere utilizzate per
i \emph{buffer}, ovvero delle aree dove NIC può immagazzinare i dati da inviare
o da ricevere.

Al momento dell'invio di un datagramma, esso verà salvato in un buffer, quindi
inviato in rete appena possibile, quindi cancellato dalla memoria.

\subsection{Ricezione di dati}
Per la ricezione di datagrammi dei buffer devono essere sempre disponibili,
cosicchè NIC possa utilizzarli una volta ricevuto il pacchetto.

Il datagramma ricevuto sarà inserito dal DMA in un buffer, un interrupt nel
frattempo creeà nuovi buffer per possibili nuovi messaggi, quindi il sistema
legge il buffer copiandolo nel buffer dell'applicazione, direzionando l'output
verso l'utente corretto, quindi la memoria viene liberata.

Il comportamento standard di un applicazione che implementa networking 
implementa un meccanismo bloccante per cui l'applicazione rimane in attesa
fintanto che il buffer non è stato scritto al suo interno.

\subsection{Trasmissione di dati su bus ring based}
La problematica principale del DMA di tipo \emph{bus mastering} è la
necessità di dover programmare in anticipo i buffer allocati, limitandone
la capacità.

Fintanto che il DMA è impegnato in un operazione perchè in attesa del
processore verrano persi i dati nel frattempo ricevuti dalla rete.

Per risolvere questo problema si implementa una struttura dati più complicata,
integrando una piccola CPU all'interno della NIC (ring based).

\subsubsection{Ring}
Il ring è composto da un anello contenente più di un buffer, e due puntatori,
un puntatore di inizio ed uno di fine gestiti da due processori diversi;
la CPU si occupa di aggiungere elementi all'interno del ring quindi gestirà
l'indice di fine, il puntatore al primo elemento viene gestito dal processore
all'interno della NIC.
SI ottiene una coda di buffer che si possono aggiungere o togliere a seconda
se si sta gestendo un invio o ricezione di dati.

Quindi la NIC legge i buffer in memoria e usa un flag per impostare se il
buffer è libero oppure se deve essere processato.

L'unico tipo di comunicazione tra processore e DMA è la trasmissione della 
quantità di buffer che dovranno essere utilizzati al DMA.

\subsection{Commutazione}
Si rende necessario, per una questione di efficienza, un circuito fisico $R$
attraverso cui trasmettere i pacchetti.

I pacchetti sono inviati come datagrammi, e nella struttura includono il
destinatario e le informazioni trasmesse.

\subsection{Dataflow}
L' \emph{host} compila ed invia il buffer in locale al buffer di ricezione di
$R1$(router), il cui compito è di leggere il datagramma, ricavarne il
destinatario, quindi di inviarlo: l'hos può quindi cancellare il buffer.

La parte ricevente ($R2$) copia il datagramma nel suo buffer di ricezione,
quindi $R1$ può rimuovere dal proprio buffer il datagramma.
Quindi $R2$ invia il datagramma all'effettivo destinatario con un altro
\emph{store and forward}.

Ogni passaggio di trasmissione del datagramma viene chiamato \emph{hop}.

\subsection{Routing}
Identifichiamo per il processo di routing l' \emph{host}, gli endpoint della
trasmissione di pacchetti e i \emph{router}, calcolatori.
Più router possono essere collegati tra loro.

\section{Internet protocol stack}
\begin{enumerate}
  \item Physical
  \item Datalink
  \item Network
  \item Transport
  \item Application
\end{enumerate}

\subsection{Protocollo fisico}
L'unico protocollo fisico sopravissuto ed in uso è il protocollo ethernet.

Ogni richiesta sarà composta da: header, payload e trailer.

Il protocollo definisce che ad ogni host venga assegnato un indirizzo
\emph{MAC}, composto da $6byte$.

L'header conterrà come dati gli indirizzi IP e le porte dei nodi comunicanti.

Il trailer è necessario per il controllo di integrità nella rete ethernet, l'
algoritmo utilizzato è \textbf{CRC32}.

\emph{Il datalink utilizza le stesse modilità di funzionamento del protocollo
fisico}.

\subsection{Network protocol}
Il protocollo network utilizza come gestore degli indirizzi i protocolli
\emph{IPv4 e IPv6}.
Ipv4 utilizza $4byte$ per il funzionamento, significa che può gestire fino a 
circa $4$ miliardi di indirizzi, mentre ipv6 ne può utilizzare molti di più.

I primi sviluppatori dei protocolli non avevano la possibilità
di interagire direttamente con i sistemi operativi, motivo per cui si è dovuto
introdurre il concetto di \textbf{porta di comunicazione}, un numero intero
in $16bit$, con cui si identifica univocamente un numero di una porta ad una
tipologia di applicazione.

Le porte da $0$ a $1023$ sono state dedicate ai protocolli internet ed alle
applicazioni più comuni, un altra serie di porte sono sempre utilizzate per
il networking  mentre le restanti sono le così dette porte effimere che vengono
utilizzate dalle applicazioni in locale.

Quindi se la porta è la parte che un calcolatore espone ad internet, il
\emph{socket} è il corrispettivo componente nel nodo in locale, così da
associare alla richiesta una macchina ed applicazione in particolare.


Sono usati i protocolli \emph{UDP} e \emph{text}
Il protocollo \emph{UDP} invia datagrammi mentre il protocollo \emph{TCP}
permette di implementare una modalità di comunicazione più avanzata, di tipo
\emph{stream}, che permette di interagire con l'interfaccia di comunicazione
come se fossero file.

Gli indirizzi \textbf{IP} sono utilizzati per essere associati agli indirizzi
\textbf{MAC} nel network.

IL \emph{TTL} (Time To Live) definisce da quanto tempo un pacchetto è nel
network.

\emph{RFC} (Remote Function Call) cerca prima di chiamare le versioni dei
protocolli più recenti prima di doverne eseguire uno più vecchio.

\subsubsection{Indirizzi}
Ad ogni macchina corrisponde un IP ($IPv4$ o $IPv6$)
\begin{figure}[H]
    \centering
    \incfig{indirizzi}
    \caption{indirizzi}
    \label{fig:indirizzi}
\end{figure}

Un indirizzo $IPv4$ in default è composto da:
\begin{itemize}
  \item $29bit$ dedicati al network;
  \item $3bit$ dedicati all'host.
\end{itemize}
Per ovviare a problemi di carenza di indirizzi IP, si introduce il concetto
di \emph{NETMASK}, una parola composta da una serie di $1$ ed una di $0$, gli
$1$ rappresentano quanti bit sono dedicati al network, gli $0$ per l'host.

\paragraph{Indirizzi privati}
Sono definiti nel protocollo \emph{RFC 1918} dichiara che un indirizzo
\textbf{IPv4} deve essere composto da $4$ numeri decimali seguiti da punto,
inoltre dichiara che gli indirizzi che cominciano con $10.$ appartengono alla
\emph{classe A}, questi indirizzi sono i cosiddetti indirizzi privati,
utilizzabili solo in \textbf{LAN}.

Gli indirizzi di classe \textbf{B} cominciano da $192.168.0.0$ per arrivare
fino $192.168.255.255$, sono sempre indirizzi privati.


Si può configurare il router per usare il \textbf{NAT} (Network Access
translator), con lo scopo di assegnare due indirizzi IP allo stesso router,
con lo scopo di esporne uno alla rete locale e l'altro a quella internet.

Con il NAT è possibile convertire indirizzi pubblici in indirizzi privati,
direzionando i messaggio alla macchina corretta sostituendo il proprio
indirizzo alla richiesta assegnando una porta libera.

\subsection{Trasporto}
Vengono utilizzato i protocolli \emph{UDP} e \emph{TCP} e la tecnica della
\emph{3 way handshake}.

\begin{figure}[ht]
    \centering
    \incfig{3some}
    \caption{Three way handshake}
    \label{fig:3some}
\end{figure}

Il trasporto comincia con l'invio di un messaggio \textbf{SYN} (accompagnato
da \textbf{SEQ}), all'arrivo
il server risponde con un messaggio che viene chiamato \textbf{SYN-ACK}, la
\textbf{SEQ} e la \textbf{ACK};
quindi l'host invia un uteriore messaggio chiamato \textbf{ACK}, oltre che
i dati.

L'header TCP contiene una serie di informazioni tra le quali i \textbf{FLAG}
SYN e ACK oltre che la porta sorgente (\textbf{SPORT}) e di destinazione
(\textbf{DPORT}).


Vengono utilizzate $2$ parole, la \textbf{SEQ} e la \textbf{ACK}, permettono
la di realizzare l'astrazione dello stream di byte; questi numeri per motivi
di sicurezza non partono da $0$, ma da un valore casuale; lo scopo della
prima squenza di passaggi è di scambiare queste informazioni.

Inoltre \textbf{SEQ} segnala anche l'inizio dello stream di byte da leggere.

Se il messaggio viene suddiviso in tanti datagrammi cumulativi, viene inviato un unico \textbf{ACK cumulativo}.

Il \textbf{timeout} è il tempo che il processo impiega ad inviare e ricevere i
messaggi, dipende dalla lunghezza dei messaggi e dalla velocità dei
calcolatori.

La \textbf{media mobile} dà più peso agli ultimi valori, si rende necessaria
per calcolare il tempo di timeout.
\[
  \text{Estimated }=(1-\alpha)\times\text{ Estimated}+\alpha\text{ SampleRTT}
\]

Inoltre si usa una stime della variabilità:
\[
  \text{Dev}=(1-\beta)\times\text{ Dev}+\beta\times\mid\text{SampleRTT}-
  \text{Estimated}\mid
\]

I parametri $\alpha,\beta$ hanno un valore consigliato: $\alpha=\dfrac{1}{8}$,
$\beta=\dfrac{1}{4}$.

Quindi inizialmente viene assegnato un valore fisso viene inizialmente
assegnato ad \emph{Estimated} e \emph{Dev}, quindi il primo messaggio viene
inviato, quindi nel primo scambuio, in cui il client riceve il \emph{SYN-ACK},
si ottiene il primo \emph{SampleRTT}.

Quindi si possono comparare i valori \emph{Estimated} e \emph{Dev},
(Deviazione) ogni scambio di \emph{ACK e SIN-ACK}.

Il \textbf{timeout sarà}:
\[
  \text{Time-out}=\text{Estimated}+4\times\text{ Dev}
\]

\subsubsection{Flow control}
Il motivo principale della perdita di messaggi è la mancanza di
sincronizzazione, un tempo era molto probabile ricevere messaggi corrotti,
ora la perdita di messaggi è legata alla mancanza di buffer liberi di
ricezione.

Il \textbf{mittente} avrà un buffer di ricezione (TX buffer), così come
il destinatario (RX buffer), nei quali i datagramm vengono ricevuti.

Per garantire che vi siano sempre buffer di ricezione liberi sarà necessario
comunicare all'interlocutore il numero di buffer allocati.

All'interno dell'header \textbf{TCP} vi sarà un campo \textbf{receive window}
che contiene la grandezza dei buffer di ricezione del destinatario.

I messaggi saranno inviati o alla ricezione di un \textbf{ACK} o alla scadenza
del \textbf{TIMEOUT}.

In generale, \emph{UDP} utilizza un approccio ottimistico per cui in una
comunicazione fluida è più efficiente rispetto a \emph{TCP} che si accerta
che ogni messaggio venga ricevuto correttamente.

\subsubsection{Congestion control}
Il \textbf{TCP} implementa anche un controllo di congestione oltre al \emph{
flow control}, è necessario per ridurre la perdita di messaggi sulla reta,
ma non da parte del destinatario, ma dai \emph{router} che implementano lo
\emph{store and forward}.

Quindi è il \emph{router} che rileva la congestione di dati (mancanza di
buffer) e rallenta la ricezione di dati.

Il congestion control avviene in tre fasi:
\begin{enumerate}
  \item \textbf{slow start}: invio di pochi bytes;
  \item dopo l' \textbf{ACK}, raddoppio la quantità di datagrammi inviati;
  \item quindi se i dati inviati diventano troppo grandi, si torna ad un
    passaggio indietro e si aumentano a passi i dati inviati.
\end{enumerate}

\section{Application layer}
L'ultimo livello di \emph{networking} è l'implementazione a livello
applicativo.
Per cui si rende necessaria un interfaccia per applicazione, \textbf{API}, si
userà lo standard \textbf{POSIX}; si utilizzano i \textbf{socket}.

\paragraph{File descriptor:}
un processo che permette di interagire con il \textit{file system}, permettono
di stabilire un collegamento tra i buffer a livello aplicativo e quelli di
sistema.

Il file descriptor (\textbf{fd}) creato, per esempio, alla chiamata di \textit{open},
sarà un \textbf{int} che rappresenta l'indice del file nell'array dei file
in utilizzo.

Nel caso di errore, \emph{fd} sarà un numero negativo.

\subsection{System calls sui socket di tipo datagram}
Ogni \textbf{system call} è bloccante, nel caso di un errore il programma si
ferma.
Si possono riprogrammare le \emph{syscall} per renderle non bloccanti, con la
gestione del messaggio di errore \textbf{E\_WOULDBLOCK}.

\subsubsection{Socket system call}
Quindi alla \textit{syscall} su \textbf{socket}, verrà creato un nuovo \emph{
fd}.

\begin{lstlisting}[caption=Esempio di syscall a socket]
int fd = socket(int domain, int type, int protocol);
\end{lstlisting}

Dove:
\begin{itemize}
  \item \textbf{domain:}
    \begin{itemize}
      \item \emph{AF\_UNIX}: locale, nomi di file;
      \item \emph{AF\_INET}: $IPV4$;
      \item \emph{AF\_INET6}: $IPV6$;
    \end{itemize}
  \item \textbf{type:}
    \begin{itemize}
      \item \emph{SOCK\_DGRAM}: $IPv4$;
      \item \emph{SOCK\_STREAM}: $IPv6$;
      \item \emph{SOCK\_RAW}: utilizzato per saltare i passaggi di trasporto e
        di rete fino a \emph{datalink};
    \end{itemize}
  \item \textbf{protocol:} esiste uno ed uno solo per type, è possibile
    lasciare il predefinito ($0$).
\end{itemize}

\subsubsection{Bind system call}
La \emph{syscall} \textbf{bind} è necessaria per associare un indirizzo ad un
\textbf{socket}.

\begin{lstlisting}[caption=Esempio di syscall a bind]
int fd = bind(int fd, const struct sockaddr*addr, socklen_t addrlen);
\end{lstlisting}

Dove:
\begin{itemize}
  \item \textbf{fd:} il file descriptor del \emph{socket} riferimento;
  \item \textbf{addr:} alloca lo spazio necessario per salvare il numero
    di $bit$ corretti;
  \item \textbf{addrlen:} la lunghezza dell'indirizzo da salvare (o del nome
    del file).
\end{itemize}

\subsubsection{Send to system call}
\begin{lstlisting}[caption=Esempio di syscall sendto]
ssize_t sendto(int fdA, const void*buf, size_t len, int flags, const struct sockaddr*addr, socklen_t addrlen);
\end{lstlisting}

Dove:
\begin{itemize}
  \item \textbf{fdA:} file descriptor del socket da cui inviare;
  \item \textbf{buf:} buffer da cui inviare;
  \item \textbf{len:} lunghezza del messaggio;
  \item \textbf{flags:} ;
  \item \textbf{sockaddr:} indirizzo del socket (comprensivo di porta);
  \item \textbf{socklen:} lunghezza del socket;
\end{itemize}

\subsubsection{Receive from system call}
\begin{lstlisting}[caption=Esempio di syscall receivefrom]
ssize_t recvfrom(int fdB, const void*buf, size_t len, int flags, struct sockaddr*addr, socklen_t addrlen);
\end{lstlisting}

Dove:
\begin{itemize}
  \item \textbf{fdA:} file descriptor del socket in cui ricevere;
  \item \textbf{buf:} buffer nel quale scrivere il messaggio;
  \item \textbf{len:} lunghezza del messaggio;
  \item \textbf{flags:} ;
  \item \textbf{sockaddr:} indirizzo del socket (comprensivo di porta del
    sender);
  \item \textbf{socklen:} lunghezza del socket;
\end{itemize}

\subsubsection{Connect system call}
\begin{lstlisting}[caption=Esempio di syscall connect]
int recvfrom(int fdA, const struct sockaddr*addr, socklen_t len);
\end{lstlisting}

\subsubsection{Listen system call}
\begin{lstlisting}[caption=Esempio di syscall listen]
int recvfrom(int fdB, int backlog);
\end{lstlisting}

\subsubsection{Accept system call}
\begin{lstlisting}[caption=Esempio di syscall accept]
int accept(int fdB, 2 parametri);
\end{lstlisting}
Il ritorno della funzione è un \textbf{fd}.
Per ogni successo il server associa un nuovo socket.

\subsection{Esempio di scambio di datagrammi}
Il server ed il client creano un socket di tipo \emph{datagram}, il server
esegue un \emph{bind}.
Il client invece esegue una \emph{syscall} di tipo \textbf{sendto}, il server
utilizza \textbf{recvfrom}, etc..

\subsection{Esempio di scambio di stream}
Il server ed il client creano socket di tipo \emph{stream}, il server esegue
la \emph{bind}, il client la \emph{connect}, per ogni \emph{send} e
\emph{receive} non è più necessario lo scambio di indirizzo.
Il server deve attivare \textbf{listen} prima dell'arrivo della \emph{connect}.


\chapter{Introduzione ai sistemi operativi}
\begin{figure}[ht]
  \centering
  \incfig{struttura}
  \caption{Struttura base di un calcolatore}\label{fig:struttura}
\end{figure}

Lo scopo di questa sezione è di virtualizzare le componenti fisiche di un
calcolatore per renderne più facile l'utilizzo.

Sinonimi di \emph{sistema operativo} sono \textbf{macchina virtuale} (virtual
machine) e \textbf{gestore dlle risorse} (resource manager).

\section{Virtualizzazione della CPU}
L'unica nota in questo programma è l'utilizzo della funzione \emph{spin ()}, che
controlla l'orologio di sistema fino a quanto $1$ minuto non è passato, quindi
riprende l'esecuzione.

\lstinputlisting[caption=Interazione con la  CPU usando la funzione spin()]{listings/cpu.c}

Da notare che se il programma viene eseguito con un comando del tipo:
\begin{lstlisting}[language=bash,caption=Esecuzione in concorrenza]
prompt> ./cpu A & ./cpu B &./cpu C & ./cpu D
[7353]
[7354]
[7355]
[7356]
A
B
D
C
A
B
D
C
A
...
\end{lstlisting}

Possiamo osservare che sembra che vengano eseguiti simultaneamente quattro
programmi diversi, ciò è dovuto alla \emph{virtualizzazione della CPU}.

\section{Memoria virtualizzata}
\lstinputlisting[caption=Dimostrazione dell'utilizzo di indirizzi virtuali]{listings/mem.c}

La memoria \textbf{RAM} è generalmente virtualizzata, difatti se eseguiamo
il programma precedente con un comando tipo:
\begin{lstlisting}[language=bash,caption=Esecuzione in concorrenza]
prompt > ./mem &; ./mem &
[1] 24113
[2] 24114
(24113) address pointed by p: 0x200000
(24114) address pointed by p: 0x200000
(24113) p: 1
(24114) p: 1
(24113) p: 2
(24114) p: 2
(24113) p: 3
(24114) p: 3
(24113) p: 4
(24114) p: 4
...
\end{lstlisting}

Possiamo osservare il funzionamento della virtualizzazione della memoria,
infatti ogni programma è lanciato in un proprio \textbf{virtual address space};
semplifica la gestione della memoria per il programma ed è mappato ad una
sezione della memoria fisica.

\section{Periferiche}
Nell'interazione di un calcolatore con le periferiche, è il kernel che
interagisce direttamente con l' hardware, mentre le applicazioni possono
utilizzare le periferiche solo attraverso il kernel.


Vi sono $4$ diversi livelli di privilegi per interagire con i dispositivi:
\begin{itemize}
  \item $0$, il livello più privilegiato, utilizzato dal kernel;
  \item $1$ solitamente non utilizzato;
  \item $2$ solitamente non utilizzato;
  \item $3$ il livello meno provilegiato, utilizzato da tutti gli utenti
    compreso l'utente di \emph{root}.
\end{itemize}

Se l'utende richiede un interazione con una periferica, un comando di
\emph{TRAP} viene inviato al kernel, il quale restituise un comando di
\emph{return from TRAP}.

\section{Memoria}
\subsection{Spazi di indirizzamento}
Ad ogni processo in esecuzione corrisponde un \emph{address space}, una
virtualizzazione della memoria che comporta diversi vantaggi, tra cui non
permettere ai programmi di modificare memoria non nella propria \emph{address
space} ed in generale un efficienza maggiore.

Lo spazio di indirizzamento non può essere più grande della memoria fisica ed
occupa una sezione di memoria contigua.

Utilizzando questo metodo abbiamo due problemi di frammentazionem interna ed
esterna.


Un processo è composto da:
\begin{itemize}
  \item \textbf{codice:} statico, regione fissa;
  \item \textbf{dati:} regione fissa;
  \item \textbf{stack:} regione dinamica, può crescere o diminuire;
  \item \textbf{heap:} regione dinamica.
\end{itemize}


Quindi l'address space è composto, in ordine, da:
\begin{itemize}
  \item \textbf{codice del programma:} dove risiedono le istruzioni;
  \item \textbf{heap:} contiene i dati creati con i \emph{malloc} e le
    strutture dinamiche, cresce verso il basso;
  \item \textbf{memoria libera:} che può essere occupata dallo heap e dallo
    stack;
  \item \textbf{stack:} dove risiedono le variabili locali, i valori di ritorno
    delle funzioni, etc. e cresce verso l'alto.
\end{itemize}


All'interno del sistema operativo è contenuta una struttura dati che contiene
informazioni relative ai processi in memoria, \textbf{PCB} (Process Control
Bloc).

\subsection{Virtualizzazione}
Gli obiettivi della virtualizzazione sono: trasparenza, efficienza e protezione
ovvero un programma non deve accorgersi che sta utilizzando solo una porzione
della memoria disponibile.


Per implementare queste funzioni si utilizza una \textbf{MMU}; al suo interno
salviamo un registro chiamato base, il quale cambia a seconda del programma
in esecuzione.

Questa implementazione non è corretta in quanto non abbiamo indicato un limite
al valore che possiamo assegnare alla base, quindi i processi potrebbero
accedere a memoria successiva alla loro, per cui sono gestiti anche i limiti
su quali valori posso assegnare al registro base.

Nel caso di un \textit{segmentation fault} (errore relativo all'accesso di
memoria non propria) vengono sollevati gli \textbf{interrupt handlers}, e gli
\textbf{exception handlers}, i quali gestiscono le eccezioni.
Altri errori simili che sollevano un interrupt sono la divisione per $0$.

Gli interrupt possono anche essere utilizzati per interagire con il sistema
operativo, il sistema operativo può disabilitare gli interrupt.

\subsection{Frammentazione}
La frammentazione interna è composta dalla memoria allocata per un programma,
ma non utilizzata.

La frammentazione esterna invece è la memoria che non viene utilizzata perchè
non abbastanza grande da ospitare un processo e tra due address spaces.

Il \textit{compattamento} della memoria non è una soluzione attuabile in quanto
sarebbe estrememante inefficiente.

Esistono diversi metodi per poter utilizzare le porzioni di memoria
frammentata.

\subsubsection{Segmentazione}
La segmentazione è utilizzata per poter utilizzare la memoria frammentata
esternamente.
Si utilizza la \textit{rilocazione dinamica}, ovvero la traduzione di indirizzi
logici in indirizzi fisici attraverso una funzione di rilocazione implementata
nell'hardware, \textbf{MMU}.


\paragraph{Registri base limite}
Nell' $8086$ in $16bit$ la segmentazione era classificata in:
\begin{itemize}
  \item \textbf{CS:} code segment;
  \item \textbf{DS:} data segment;
  \item \textbf{SS:} stack segment;
  \item \textbf{ES:} extra segment.
\end{itemize}

In questa versione, la \textbf{MMU} mantiene una lista di registri \textit{
base/limite}, ad ogni coppia corrisponde un segmento di memoria frammentata,
la base contiene il primo indirizzo di memoria utile, mentre il limite quanta
memoria contigua è disponibile.

Nel momento in cui il processore genere un indirizzo virtuale prende in
considerazione questi registri; l'indirizzo viene generato in fase di \textit{
fetch} dell'istruzione.
Se l'indirizzo generato supera il valore contenuto nel registro limite, un
eccezione viene generata in quanto l'indirizzo richiesto non apparterrebbe al
segmento.

\paragraph{Tabelle dei descrittori}
L'architettura $386$, risolve due problemi relativi alla precedente soluzione:
una più veloce risoluzione dei permessi ed il limite di registri salvabili
in MMU.


Per ottenere l'indirizzo di memoria giusto viene ora utilizzata una parola che
indica sia il \emph{segmento} di riferimento che l'offset all'interno del
segmento (indirizzo virtuale).


Inoltre viene introdotta una struttura dai tabellare, mantenuta all'interno
della memoria riservata del sistema operativo, detta \textit{tabella dei
descrittori}, contenente i dati relativi alla coppia dei registri \textit{base/
limite} di ogni segmento, il bit di protezione e la gestione dei privilegi
DPL).
La MMU conterrà soltanto due registri, \textbf{STLR} (Segment Table Limit
Register) la dimensione della tabella dei segmenti e \textbf{STBR} (Segment
Table Base Register) l'indirizzo base della tabella dei segmenti.

\begin{itemize}
  \item \textbf{CS:} code segment, con una local descriptor table;
  \item \textbf{DS:} data segment, con una global descriptor table;
  \item \textbf{SS:} stack segment;
  \item \textbf{ES:} extra segment.
\end{itemize}
\[
  MAX(CPL \quad RPL) \leq DPL \quad BASE + OFFSET
\]

I segmenti di codice (\textbf{CS}) possono essere utilizzati da più processi
uguali.

Quindi le funzioni \textit{Jmp, Call e Ret} per accedere ad un segmento
utilizzano due dati: \textbf{offset} e \textbf{selettore offset}.

Definiamo \textbf{EIP} (Instruction Pointer) e \textbf{ESP} (Stack Pointer),
rispettivamente appartengono ad un determinato \textbf{CS} (Code Segment) e
\textbf{SS} (Stack Segment).

\textbf{EIP} e \textbf{CS} alla chiamata di \textit{Jmp, Call o Ret}, verranno
salvati nello stack (\textbf{ESP}).



Il \textbf{Call Gate} è uno strato che si aggiunge all'indirizzo virtuale per
non permettere l'erroneo accesso a sezioni di memoria non utilizzabili dall'
utente o dal programma.


Ogni livello di privilegio ha uno \textbf{stack separato}, per offrire una
maggiore stabilità del sistema, quindi se una \emph{Call} cambia il livello
di privilegio dovrà anche cambiare lo \emph{stack}: se il livello di privilegio
si riduce, tutti i registri che contengono riferimenti a sezioni di memoria
privilegiata vengono azzerati.

\paragraph{Interrupt}
Se l'utente è kernel, al momento dell'arrivo di un interrupt non c'è bisogno
di cambiare lo stack, altrimenti sì.

All'arrivo di un interrupt \textbf{eflag} (flag dell'interrupt) viene salvato,
quindi l'interrupt viene disattivato, \textbf{iret} esce dall'interrupt
handler.


\subsection{Paginazione}
Per il processo la paginazione è trasparente.
La memoria fisica è divisa equamente in \textbf{page frame}, la
virtualizzazione corrispondente sono le \textbf{pagine}, ovvero le pagine sono
ordinate, mentre la rappresentazione nel \textbf{page frame} non è
necessariamente ordinata, ma i \textbf{page frame} sono numerati.
Le pagine sono grandi $4096byte$, in generale sempre potenze di $2$.


Quindi si utilizza una tabella di appoggio, lunga tanto quanto il numero di
pagine, per associare \emph{page frames} a \emph{pagine}, in questa tabella la
colonna \textbf{present}(flag), vero se la pagina è utilizzata nella memoria
fisica.
Una tabella simile è creata ed associata per ogni processo, la tabella è
creabile e modificabile solo dal kernel.


Dato un indirizzo virtuale si divide per la grandezza della pagina per ottenere
il numero della pagina, quindi dalla tabella di appoggio si seleziona il \emph{
page frame} corrispondente, quindi si utilizza il resto della divisione per
determinare l'indirizzo di \textbf{memoria fisica} corrispondente.
\[
  \text{Indirizzo fisico} = \text{Indirizzo virtuale}/3 * \text{Grandezza delle
  pagine} + \text{Indirizzo virtuale} \% 3 
\]

Stabilita la grandezza della pagina si interpretano i primi $bit$ dell' 
indirizzo virtuale per trovare il numero di pagina, mentre il restante si usa
come \emph{offset}, dato l'indirizzo della base si appende l'offset per
ottenere l'indirizzo fisico.


La paginazione usa appunto le \textbf{paging tables}, per ogni riga abbiamo:
\begin{itemize}
  \item il \textbf{PPN} (Physical Page Number);
  \item un bit di validità;
  \item un bit di presenza (se la pagina è presente);
  \item un bit di accesso (se ha già acceduto);
  \item un bit di modifica;
  \item dei bit per il livello di privilegio (r,w,x).
\end{itemize}


Per migliorare l'efficienza si utilizza una cache, la \textbf{TLB}
(Translation Lookaside Buffer), da un \textbf{VPN} (Virtual Page Number),
si ottiene il \textbf{PPN} (Physical Page Number) ed il \textbf{PFN}
(Physical Frame Number).


Per migliorare ulteriormente la paginazione, non si utilizzano i bit di
accesso e modifica, sarà quindi il kernel a tenere traccia di questi valori,
lo stesso ragionamento si può fare per i bit sui permessi.

\paragraph{Osservazione}
Per indiciare una ram di $4GB$ sono necessari $32bit$, un indirizzo virtuale
è composto da $20bit$ per selezionare la pagina e:
\begin{itemize}
  \item $v=1bit$;
  \item $p=1bit$;
  \item $r/w/x=3bit$;
  \item $d=1bit$.
\end{itemize}

Quindi abbiamo un totale di $27bit$ ($32bit$ normalizzati), per una RAM da
$4GB$ la \emph{paging table} occuperebbe $400mb$ di memoria contigua nella
RAM.

Per rendere più efficiente la tabella, i $20bit$ di indirizzamento di dividono
logicamente in 2 blocchi da $10bit$, il primo per indirizzare il blocco di
memoria, il secondo per ottenere l'effettivo indirizzo fisico.

Il primo blocco in realtà indirizza su di una nuova tabella di appoggio il cui
contenuto sarà se la pagina è occupata o no e l'indirizzo della tabella
corrispondente, quindi si usa la tabella di appoggio ed il secondo blocco di
$10bit$ per ottenere l'indirizzo fisico.

\subsubsection{Base limite, Segmentazione e paginazione}
Compariamo i sistemi di suddivisione della memoria: \textbf{B+L}, \textbf{
segmentazione} e \textbf{paginazione}.

B+L:
\begin{itemize}
  \item $?$ spazio di indirizzamento contiguo in \emph{RAM};
  \item $-$ frammentazione interna;
  \item $-$ frammentazione esterna;
  \item $-$ condivisione codice;
  \item $+$ isolamento\\protezione;
  \item $-$ efficiente;
  \item $-$ spazio di indirizzamento contiguo.
\end{itemize}

Segmentazione:
\begin{itemize}
  \item $?$ spazio di indirizzamento contiguo nel segmento, ma non nel processo
    ;
  \item $-$ frammentazione interna;
  \item $-$ frammentazione esterna;
  \item $\mp$ condivisione codice;
  \item $+$ isolamento\\protezione;
  \item $\pm$ efficiente;
  \item $-$ spazio di indirizzamento contiguo.
\end{itemize}

Paginazione:
\begin{itemize}
  \item $-$ spazio di indirizzamento non contiguo;
  \item $\pm$ frammentazione interna ($4k$);
  \item $-$ frammentazione esterna;
  \item $\mp$ condivisione codice;
  \item $+$ isolamento\\protezione;
  \item $\pm$ efficiente;
  \item $-$ spazio di indirizzamento contiguo.
\end{itemize}

\section{System calls}
\subsection{fork}
La funzione \emph{fork} crea un processo quasi identico al processo padre e
restituisce un numero positivo chiamato \emph{pid\_t}.
Ogni processo ha un \emph{ppid} ( \emph{pid} del padre).

\subsection{exit}
La funzione \emph{exit} ritorna $0$ se il processo si chiude correttamente,
altrimenti un numero.

\subsection{wait}
Ritorna un \emph{pid\_t} quando il processo passato come argomento termina. 

\subsection{exec}
Sostituisce il processo corrente con uno nuovo, l'argomento è un PATH.

\subsection{open}
Apre un file e gli assegna il primo file descriptor disponibile.

\subsection{dup}
Duplica un fd creandone uno nuovo che punta allo stesso oggetto.

\subsection{dup2}
La dup2 ha due parametri, il primo il vecchio file descriptor ed il secondo il
nuovo, duplica il primo parametro nel secondo, fa una close del secondo
parametro nel caso fosse già aperto.

\subsection{pipe}
\begin{lstlisting}[
  language=bash,
  escapeinside={(£}{£)},
  caption={
    Esempio di pipe (|)
  }
]
# (1)
ls -l | grep cammello

# (2)
ls -l > camelus
grep cammello < camelus
\end{lstlisting}

L'operatore $|$ permette di chiamare la syscall \textit{pipe()} la quale
permette di eseguire l'operazione esclusivamente in memoria, difatti nel
secondo esempio \textit{camelus} sarebbe scritto su disco e, su risultati
molto grandi sarebbe molto più lento ed inefficiente.

Quindi \textit{pipe(fds)} prende come argomento un array di file descriptors
ed utilizza un buffer riservato al kernel per svolgere la lettura.

Il funzionamento per il  comando $(1)$ sarà:
\begin{itemize}
  \item il processo padre chiama la \textit{pipe()}, $2$ \textit{fork()},
    \textit{close(3)}, \textit{close(4)} e $2$ \textit{wait()};
  \item la prima figlia \textit{ls -l} chiama:
    \begin{itemize}
      \item \textit{close(1)};
      \item \textit{dup(3)};
      \item \textit{close(2)};
      \item \textit{close(4)} dove $4$ è il fd della pipe;
      \item \textit{exec(...)};
    \end{itemize}
  \item la seconda figlia \textit{grep cammello} chiama:
    \begin{itemize}
      \item \textit{close(0)};
      \item \textit{dup(4)}, dove $3$ è il file descriptor della fork;
      \item \textit{close(2)};
      \item \textit{close(4)} dove $4$ è il fd della pipe;
      \item \textit{exec(...)};
    \end{itemize}
\end{itemize}

La shell è il programma che attende la scrittura e l'invio dei comandi, una
delle funzionalità è di espandere le variabili passate (es. *.c).

\section{Scheduling}
Lo scheduler è un programma che implementa un algoritmo di scheduling il quale,
dato un insieme di richieste di accesso ad una risorsa, stabilisce un
ordinamento temporale per l'esecuzione di tali richieste, privilegiando quelle
che rispettano determinati parametri secondo una certa politica di scheduling,
in modo da ottimizzare l'accesso a tale risorsa e consentire così
l'espletamento del servizio/istruzione o processo desiderato.

Funziona tramite i meccanismi di \emph{LDE} (user, kernel) e di context
switching.

Ricapitolando, il ciclo di vita di un processo:
\begin{itemize}
  \item init;
  \item \emph{EMBRYO};
  \item fork;
  \item ready;
  \item running;
  \item zombie.
\end{itemize}

\subsection{Turn-around Time}
\[
  T_{TA}=T_{ \text{Completion}}-T_{ \text{Arrival}}
\]

%%% This defines the bibliography file (main.bib) and the bibliography style.
%% If you want to create a bibliography file by hand, change the contents of
%% this file to a `thebibliography' environment.  For more information 
%% see section 4.3 of the LaTeX manual.
\begin{singlespace}
\bibliography{main}
\bibliographystyle{plain}
\end{singlespace}

\end{document}

